% ============================================================
%  Chapter 2 -- Compact Sets, 2.31-2.42
% ============================================================
\rsection{Compact Sets}

\begin{signpost}[The centre of the chapter]
Everything so far has been vocabulary. This section introduces the one property
that does real work, and it arrives in a form that looks arbitrary: a set is
compact if every open cover has a finite subcover.

\smallskip
\textbf{Why that definition.} Compactness is the hypothesis that converts
``for each point there is some local fact'' into ``there is one fact for the
whole set.'' Local statements come with a neighbourhood attached; a finite
subcover lets you take a minimum over finitely many of them, and a minimum of
finitely many positive numbers is positive (the move already used at 2.20 and
2.24(c)). Every application in Chapter 4 --- boundedness 4.15, attainment 4.16,
uniform continuity 4.19 --- is that one manoeuvre.

\smallskip
\textbf{The plan of the section.} 2.33--2.37 develop compactness abstractly, in
any metric space. 2.38--2.40 then do the hard work of exhibiting compact sets
in $\R^k$, culminating in ``every $k$-cell is compact.'' 2.41 collects the
payoff: in $\R^k$, compact means closed and bounded. Note the division ---
2.33--2.37 are general, 2.38--2.42 are about $\R^k$ and are \emph{false} in
general metric spaces.
\end{signpost}

\ritem{2.31}{Definition}
By an \emph{open cover} of a set $E$ in a metric space $X$ we mean a collection
$\{G_\alpha\}$ of open subsets of $X$ such that $E \subset \bigcup_\alpha
G_\alpha$.

\ritem{2.32}{Definition}
A subset $K$ of a metric space $X$ is said to be \emph{compact} if every open
cover of $K$ contains a finite subcover.

More explicitly, this requirement is that if $\{G_\alpha\}$ is an open cover of
$K$, then there are finitely many indices $\alpha_1,\dots,\alpha_n$ such that
\[ K \subset G_{\alpha_1} \cup \cdots \cup G_{\alpha_n}. \]

The notion of compactness is of great importance in analysis, especially in
connection with continuity (Chap.\ 4).\kw{Read the quantifiers carefully.
\emph{Every} cover must admit a finite subcover --- exhibiting one cover that
happens to have a finite subcover proves nothing. To show a set is \emph{not}
compact it suffices to produce a single bad cover, which is why
non-compactness is usually the easy direction.}

It is clear that every finite set is compact. The existence of a large class of
infinite compact sets in $\R^k$ will follow from Theorem 2.41.

\begin{unpack}[What a cover looks like, and what fails]
Two examples settle the intuition. Both are covers of a subset of the line by
open intervals; one collapses to finitely many, the other cannot.

\medskip
\begin{center}
\begin{tikzpicture}[x=1mm, y=1mm]
  % ---- compact: [0,1]
  \draw[axis] (-3,14) -- (62,14);
  \draw[seg] (4,14) -- (48,14);
  \node[ptfill, minimum size=2.6pt] at (4,14) {};
  \node[ptfill, minimum size=2.6pt] at (48,14) {};
  \node[mlbl, anchor=north] at (4,12.4) {$0$};
  \node[mlbl, anchor=north] at (48,12.4) {$1$};
  \foreach \a/\b/\h in {1/16/18, 12/30/21.5, 26/42/18, 38/52/21.5}{
    \draw[thin] (\a,\h) -- (\b,\h);
    \node[ptopen, minimum size=2.6pt] at (\a,\h) {};
    \node[ptopen, minimum size=2.6pt] at (\b,\h) {};}
  \node[lbl, anchor=west] at (66,19) {$[0,1]$: finitely many};
  \node[lbl, anchor=west] at (66,15) {always suffice};

  % ---- not compact: (0,1]
  \draw[axis] (-3,-14) -- (62,-14);
  \draw[seg] (4,-14) -- (48,-14);
  \node[ptopen, minimum size=2.6pt] at (4,-14) {};
  \node[ptfill, minimum size=2.6pt] at (48,-14) {};
  \node[mlbl, anchor=north] at (4,-15.6) {$0$};
  \node[mlbl, anchor=north] at (48,-15.6) {$1$};
  \foreach \a/\h in {26/-10, 15/-6.5, 9.5/-3, 6.7/0.5}{
    \draw[thin] (\a,\h) -- (52,\h);
    \node[ptopen, minimum size=2.6pt] at (\a,\h) {};}
  \node[lbl, anchor=west] at (66,-4) {$(0,1]$ covered by};
  \node[lbl, anchor=west] at (66,-8) {$(1/n,\,2)$: no finite};
  \node[lbl, anchor=west] at (66,-12) {subfamily reaches $0$};
\end{tikzpicture}
\end{center}
\figcap{Compactness fails for $(0,1]$ because the cover creeps toward a missing
endpoint. Any finite subfamily stops short.}

\smallskip
The failure is instructive: $(0,1]$ is bounded but not closed, and the cover
exploits exactly the missing point. Theorem 2.41 will say that in $\R^k$ these
are the only two ways compactness can fail.
\end{unpack}

We observed earlier (in Remark 2.29) that if $E \subset Y \subset X$, then $E$
may be open relative to $Y$ without being open relative to $X$. The property of
being open thus depends on the space in which $E$ is embedded. The same is true
of the property of being closed.

Compactness, however, behaves better, as we shall now see. To formulate the
next theorem, let us say, temporarily, that $K$ is \emph{compact relative to
$X$} if the requirements of Definition 2.32 are met.

\ritem{2.33}{Theorem}
\emph{Suppose $K \subset Y \subset X$. Then $K$ is compact relative to $X$ if
and only if $K$ is compact relative to $Y$.}

By virtue of this theorem we are able, in many situations, to regard compact
sets as metric spaces in their own right, without paying any attention to any
embedding space. In particular, although it makes little sense to talk of open
spaces, or of closed spaces (every metric space $X$ is an open subset of
itself, and is a closed subset of itself), it does make sense to talk of
\emph{compact metric spaces}.\kn{This is why compactness is the better notion.
``Open'' and ``closed'' are relative and must always be qualified; compactness
is \emph{intrinsic} --- a property of $K$ with its own metric, independent of
where $K$ happens to sit. That is what makes it usable as a hypothesis in
Chapter 4, where the ambient space changes from line to line.}

\begin{proof}
Suppose $K$ is compact relative to $X$, and let $\{V_\alpha\}$ be a collection
of sets, open relative to $Y$, such that $K \subset \bigcup_\alpha V_\alpha$.
By Theorem 2.30, there are sets $G_\alpha$, open relative to $X$, such that
$V_\alpha = Y \cap G_\alpha$, for all $\alpha$; and since $K$ is compact
relative to $X$, we have
\begin{equation}
  K \subset G_{\alpha_1} \cup \cdots \cup G_{\alpha_n} \tag{2.22}
\end{equation}
for some choice of finitely many indices $\alpha_1,\dots,\alpha_n$. Since
$K \subset Y$, (2.22) implies
\begin{equation}
  K \subset V_{\alpha_1} \cup \cdots \cup V_{\alpha_n}. \tag{2.23}
\end{equation}
This proves that $K$ is compact relative to $Y$.

Conversely, suppose $K$ is compact relative to $Y$, let $\{G_\alpha\}$ be a
collection of open subsets of $X$ which covers $K$, and put
$V_\alpha = Y \cap G_\alpha$. Then (2.23) will hold for some choice of
$\alpha_1,\dots,\alpha_n$; since $V_\alpha \subset G_\alpha$, (2.23) implies
(2.22).

This completes the proof.
\end{proof}

\ritem{2.34}{Theorem}
\emph{Compact subsets of metric spaces are closed.}

\begin{proof}
Let $K$ be a compact subset of a metric space $X$. We shall prove that the
complement of $K$ is an open subset of $X$.

Suppose $p \in X$, $p \notin K$. If $q \in K$, let $V_q$ and $W_q$ be
neighborhoods of $p$ and $q$, respectively, of radius less than
$\frac12 d(p,q)$ [see Definition 2.18(a)]. Since $K$ is compact, there are
finitely many points $q_1,\dots,q_n$ in $K$ such that
\[ K \subset W_{q_1} \cup \cdots \cup W_{q_n} = W. \]
If $V = V_{q_1} \cap \cdots \cap V_{q_n}$, then $V$ is a neighborhood of $p$
which does not intersect $W$. Hence $V \subset K^c$, so that $p$ is an interior
point of $K^c$. The theorem follows.\kn{Here is the manoeuvre the signpost
promised, in its purest form. For each $q$ we get a \emph{local} separation ---
two small balls, radius depending on $q$. Compactness reduces the infinitely
many $q$ to finitely many, and only then may we intersect the $V_q$: an
infinite intersection of neighbourhoods need not be a neighbourhood, exactly as
2.25 showed.}
\end{proof}

\begin{center}
\begin{tikzpicture}[x=1mm, y=1mm]
  \fill[black!8] plot[smooth cycle, tension=0.8]
    coordinates {(34,4) (56,2) (66,16) (52,26) (34,22)};
  \draw[thin] plot[smooth cycle, tension=0.8]
    coordinates {(34,4) (56,2) (66,16) (52,26) (34,22)};
  \node[lbl] at (50,14) {$K$};
  \node[ptfill, minimum size=2.6pt] at (6,14) {};
  \node[mlbl, anchor=east] at (4,14) {$p$};
  \foreach \pt in {(36,10),(50,5),(63,15),(50,24)}{
    \node[ptfill, minimum size=2pt] at \pt {};}
  \draw[guide] (6,14) circle (11mm);
  \node[lbl, anchor=south] at (6,26) {$V$};
  \node[lbl, anchor=north] at (44,-1) {$W$: finitely many balls covering $K$};
  \foreach \pt/\r in {(36,10)/7,(50,5)/8,(63,15)/8,(50,24)/8}{
    \draw[thin] \pt circle (\r pt);}
\end{tikzpicture}
\end{center}
\figcap{Each $q \in K$ is separated from $p$ by a pair of balls. Finitely many
$W_q$ already cover $K$, so the matching finitely many $V_q$ may be intersected
to give one ball around $p$ missing $K$.}

\ritem{2.35}{Theorem}
\emph{Closed subsets of compact sets are compact.}

\begin{proof}
Suppose $F \subset K \subset X$, $F$ is closed (relative to $X$), and $K$ is
compact. Let $\{V_\alpha\}$ be an open cover of $F$. If $F^c$ is adjoined to
$\{V_\alpha\}$, we obtain an open cover $\Omega$ of $K$.\kn{The trick of the
proof, and it is worth remembering: to convert a cover of $F$ into a cover of
the larger $K$, throw in the single open set $F^c$, which mops up everything
outside $F$. It is legitimately open precisely because $F$ is closed --- 2.23.}
Since $K$ is compact, there is a finite subcollection $\Phi$ of $\Omega$ which
covers $K$, and hence $F$. If $F^c$ is a member of $\Phi$, we may remove it
from $\Phi$ and still retain an open cover of $F$. We have thus shown that a
finite subcollection of $\{V_\alpha\}$ covers $F$.
\end{proof}

\begin{center}
\begin{tikzpicture}[x=1mm, y=1mm]
  % K, with K\F hatched (covered by the single set F^c)
  \fill[pattern=north east lines, pattern color=black!45, even odd rule]
    plot[smooth cycle, tension=0.8]
      coordinates {(0,4) (30,-2) (58,6) (52,26) (20,30) (-4,20)}
    plot[smooth cycle, tension=0.8]
      coordinates {(14,10) (34,8) (40,18) (24,22)};
  \draw[thin] plot[smooth cycle, tension=0.8]
    coordinates {(0,4) (30,-2) (58,6) (52,26) (20,30) (-4,20)};
  % F itself
  \fill[black!14] plot[smooth cycle, tension=0.8]
    coordinates {(14,10) (34,8) (40,18) (24,22)};
  \draw[thin] plot[smooth cycle, tension=0.8]
    coordinates {(14,10) (34,8) (40,18) (24,22)};
  % the given cover, over F only
  \foreach \pt in {(18,12),(28,10),(35,16),(22,19)}{
    \draw[thin] \pt circle (6.5pt);}
  \node[mlbl] at (27,15) {$F$};
  % labels
  \node[lbl, anchor=west] at (64,24) {hatched: $F^c$, one open set,};
  \node[lbl, anchor=west] at (64,20.5) {swallows all of $K$ outside $F$};
  \node[lbl, anchor=west] at (64,12) {circles: the given cover};
  \node[lbl, anchor=west] at (64,8.5) {$\{V_\alpha\}$, needed only on $F$};
  \node[lbl, anchor=north] at (10,-4) {$K$ compact};
  \draw[thin] (10,-2.5) -- (13,1.5);
\end{tikzpicture}
\end{center}
\figcap{Adjoining the single open set $F^c$ turns a cover of $F$ into a cover
of $K$; compactness of $K$ then hands back finitely many sets, and discarding
$F^c$ leaves a finite cover of $F$.}

\ritem{}{Corollary}
\emph{If $F$ is closed and $K$ is compact, then $F \cap K$ is compact.}

\begin{proof}
Theorems 2.24(b) and 2.34 show that $F \cap K$ is closed; since
$F \cap K \subset K$, Theorem 2.35 shows that $F \cap K$ is compact.
\end{proof}

\ritem{2.36}{Theorem}
\emph{If $\{K_\alpha\}$ is a collection of compact subsets of a metric space
$X$ such that the intersection of every finite subcollection of $\{K_\alpha\}$
is nonempty, then $\bigcap K_\alpha$ is nonempty.}

\begin{proof}
Fix a member $K_1$ of $\{K_\alpha\}$ and put $G_\alpha = K_\alpha^c$. Assume
that no point of $K_1$ belongs to every $K_\alpha$. Then the sets $G_\alpha$
form an open cover of $K_1$; since $K_1$ is compact, there are finitely many
indices $\alpha_1,\dots,\alpha_n$ such that
$K_1 \subset G_{\alpha_1} \cup \cdots \cup G_{\alpha_n}$. But this means that
\[ K_1 \cap K_{\alpha_1} \cap \cdots \cap K_{\alpha_n} \]
is empty, in contradiction to our hypothesis.
\end{proof}

\ritem{}{Corollary}
\emph{If $\{K_n\}$ is a countable collection of nonempty compact sets such that
$K_n \supset K_{n+1}$ $(n = 1,2,3,\dots)$, then $\bigcap_{n=1}^\infty K_n$ is
not empty.}

\begin{deeper}[The finite intersection property is compactness turned inside out]
2.36 is 2.32 with complements taken, and nothing more. ``Every open cover has a
finite subcover'' says: if no finite subfamily of $\{G_\alpha\}$ covers $K$,
then $\{G_\alpha\}$ does not cover $K$. Complementing every $G$ turns
``covers'' into ``has empty intersection with,'' and the statement becomes
2.36.

The corollary is the form that gets used, and it answers the worry raised at
Example 2.10(b)(iii), where nested nonempty sets had empty intersection. There
the sets were the segments $(0,x)$ --- bounded but not closed, hence not
compact. Compactness is exactly what forbids the nesting from escaping to
nothing.

Compare Theorem 2.38, the same statement for closed intervals proved directly
from the least-upper-bound property, and Yu Pin's nested interval axiom in
Chapter 1. Three routes to one fact, and Rudin uses whichever is cheapest at
the point of need.
\end{deeper}

\ritem{2.37}{Theorem}
\emph{If $E$ is an infinite subset of a compact set $K$, then $E$ has a limit
point in $K$.}

\begin{proof}
If no point of $K$ were a limit point of $E$, then each $q \in K$ would have a
neighborhood $V_q$ which contains at most one point of $E$ (namely, $q$, if
$q \in E$). It is clear that no finite subcollection of $\{V_q\}$ can cover
$E$; and the same is true of $K$, since $E \subset K$. This contradicts the
compactness of $K$.
\end{proof}

\ritem{2.38}{Theorem}
\emph{If $\{I_n\}$ is a sequence of intervals in $\R^1$, such that
$I_n \supset I_{n+1}$ $(n = 1,2,3,\dots)$, then $\bigcap_{n=1}^\infty I_n$ is
not empty.}

\begin{proof}
If $I_n = [a_n,b_n]$, let $E$ be the set of all $a_n$. Then $E$ is nonempty and
bounded above (by $b_1$). Let $x$ be the sup of $E$. If $m$ and $n$ are
positive integers, then
\[ a_n \le a_{m+n} \le b_{m+n} \le b_m, \]
so that $x \le b_m$ for each $m$. Since it is obvious that $a_m \le x$, we see
that $x \in I_m$ for $m = 1,2,3,\dots$.\kn{This is Chapter 1 doing the work
again --- the whole proof is ``take a supremum,'' which exists by 1.19. Every
$b_m$ is an upper bound for the left endpoints, so the least upper bound $x$
sits below all of them and above all the $a$'s. Note the closedness of the
intervals is essential: for $I_n=(0,1/n)$ the same construction gives $x=0$,
which lies in no $I_n$.}
\end{proof}

\begin{center}
\begin{tikzpicture}[x=1mm, y=1mm]
  \foreach \n/\a/\b in {0/2/86, 1/12/74, 2/22/64, 3/30/57, 4/36/52}{
    \pgfmathsetmacro\yy{-\n*6}
    \draw[seg] (\a,\yy) -- (\b,\yy);
    \node[ptfill, minimum size=2.4pt] at (\a,\yy) {};
    \node[ptfill, minimum size=2.4pt] at (\b,\yy) {};}
  \node[mlbl, anchor=east] at (0,0) {$I_1$};
  \node[mlbl, anchor=east] at (10,-6) {$I_2$};
  \node[mlbl, anchor=east] at (20,-12) {$I_3$};
  \draw[guide] (43,3) -- (43,-30);
  \node[ptfill, minimum size=3pt] at (43,-27) {};
  \node[lbl, anchor=north] at (43,-30) {$x=\sup a_n$ lies in every $I_n$};
  \node[lbl, anchor=west] at (90,-8) {endpoints};
  \node[lbl, anchor=west] at (90,-12) {included};
\end{tikzpicture}
\end{center}
\figcap{Left endpoints increase, right endpoints decrease, and $\sup a_n$ is
caught between. Drop the endpoints and the common point can escape.}

\ritem{2.39}{Theorem}
\emph{Let $k$ be a positive integer. If $\{I_n\}$ is a sequence of $k$-cells
such that $I_n \supset I_{n+1}$ $(n = 1,2,3,\dots)$, then
$\bigcap_{n=1}^\infty I_n$ is not empty.}

\begin{proof}
Let $I_n$ consist of all points $\mathbf{x} = (x_1,\dots,x_k)$ such that
\[ a_{n,j} \le x_j \le b_{n,j} \qquad (1 \le j \le k;\ n = 1,2,3,\dots), \]
and put $I_{n,j} = [a_{n,j},\,b_{n,j}]$. For each $j$, the collection
$\{I_{n,j}\}$ satisfies the hypotheses of Theorem 2.38. Hence there are real
numbers $x_j^*$ $(1 \le j \le k)$ such that
\[ a_{n,j} \le x_j^* \le b_{n,j} \qquad (1 \le j \le k;\ n = 1,2,3,\dots). \]
Setting $\mathbf{x}^* = (x_1^*,\dots,x_k^*)$, we see that
$\mathbf{x}^* \in I_n$ for $n = 1,2,3,\dots$. The theorem follows.\kn{The
$k$-dimensional statement is $k$ separate one-dimensional statements, because a
$k$-cell is a product of intervals and a point lies in it exactly when each
coordinate lies in the corresponding interval. This coordinatewise reduction is
used again in 2.40 and throughout Chapter 9.}
\end{proof}

\ritem{2.40}{Theorem}
\emph{Every $k$-cell is compact.}

\begin{proof}
Let $I$ be a $k$-cell, consisting of all points
$\mathbf{x} = (x_1,\dots,x_k)$ such that $a_j \le x_j \le b_j$
$(1 \le j \le k)$. Put
\[ \delta = \Bigl\{\sum_{1}^{k} (b_j-a_j)^2\Bigr\}^{1/2}. \]
Then $|\mathbf{x}-\mathbf{y}| \le \delta$, if $\mathbf{x} \in I$,
$\mathbf{y} \in I$.

Suppose, to get a contradiction, that there exists an open cover
$\{G_\alpha\}$ of $I$ which contains no finite subcover of $I$. Put
$c_j = (a_j+b_j)/2$. The intervals $[a_j,c_j]$ and $[c_j,b_j]$ then determine
$2^k$ $k$-cells $Q_i$ whose union is $I$. At least one of these sets $Q_i$,
call it $I_1$, cannot be covered by any finite subcollection of $\{G_\alpha\}$
(otherwise $I$ could be so covered). We next subdivide $I_1$ and continue the
process. We obtain a sequence $(I_n)$ with the following properties:
\begin{enumerate}[label=(\alph*), leftmargin=2.2em, itemsep=1pt, topsep=3pt]
  \item $I \supset I_1 \supset I_2 \supset I_3 \supset \cdots$;
  \item $I_n$ is not covered by any finite subcollection of $\{G_\alpha\}$;
  \item If $\mathbf{x} \in I_n$ and $\mathbf{y} \in I_n$, then
        $|\mathbf{x}-\mathbf{y}| \le 2^{-n}\delta$.
\end{enumerate}
By (a) and Theorem 2.39, there is a point $\mathbf{x}^*$ which lies in every
$I_n$. For some $\alpha$, $\mathbf{x}^* \in G_\alpha$. Since $G_\alpha$ is
open, there exists $r > 0$ such that $|\mathbf{y}-\mathbf{x}^*| < r$ implies
that $\mathbf{y} \in G_\alpha$. If $n$ is so large that $2^{-n}\delta < r$
(there is such an $n$, for otherwise $2^n \le \delta/r$ for all positive
integers $n$, which is absurd since $\R$ is archimedean), then (c) implies that
$I_n \subset G_\alpha$, which contradicts (b).

This completes the proof.
\end{proof}

\begin{center}
\begin{tikzpicture}[x=1mm, y=1mm]
  \node[rmbox, text width=29mm] (a0) at (0,22)
    {suppose: no finite\\subcover of $I$};
  \node[rmbox, text width=29mm] (a1) at (45,22)
    {halve all edges: one\\cell inherits the failure};
  \node[rmbox, text width=27mm] (a2) at (89,22)
    {nested $I_n$, diam\\$\le 2^{-n}\delta$, all bad};
  \draw[rmarrow] (a0) -- (a1);
  \draw[rmarrow] (a1) -- node[lbl, anchor=south, inner sep=2.5pt]
    {repeat} (a2);
  \node[rmbox, text width=27mm] (a3) at (89,0)
    {$\mathbf{x}^* \in$ every $I_n$};
  \draw[rmarrow] (a2) -- node[lbl, anchor=west, inner sep=3pt]
    {2.39} (a3);
  \node[rmbox, text width=31mm] (a4) at (45,0)
    {$\mathbf{x}^* \in G_\alpha$ open;\\archimedes: $I_n \subset G_\alpha$};
  \draw[rmarrow] (a3) -- (a4);
  \node[rmbox, text width=28mm] (a5) at (0,0)
    {$I_n$ covered by \emph{one}\\set --- contradiction};
  \draw[rmarrow] (a4) -- (a5);
\end{tikzpicture}
\end{center}
\figcap{The proof as a flow diagram. Completeness enters once (2.39, via
2.38), the archimedean property once, and the contradiction is that one is
finite.}

\begin{unpack}[Bisection: the shape of the argument]
This is the most important proof in the chapter, and its structure recurs
throughout analysis. Assume the conclusion fails, then manufacture a nested
sequence of ever-smaller pieces in which it keeps failing, and squeeze them
onto a single point where the failure becomes absurd.

\medskip
\begin{center}
\begin{tikzpicture}[x=1mm, y=1mm]
  % stage 0
  \draw[thin] (0,0) rectangle (26,26);
  \node[lbl, anchor=north] at (13,-2) {$I$};
  % stage 1
  \begin{scope}[xshift=34mm]
    \draw[thin] (0,0) rectangle (26,26);
    \draw[thin] (13,0) -- (13,26); \draw[thin] (0,13) -- (26,13);
    \fill[black!18] (13,0) rectangle (26,13);
    \draw[thin] (13,0) rectangle (26,13);
    \node[lbl, anchor=north] at (13,-2) {$I_1$: one quarter fails};
  \end{scope}
  % stage 2
  \begin{scope}[xshift=68mm]
    \draw[thin] (0,0) rectangle (26,26);
    \draw[thin] (13,0) -- (13,26); \draw[thin] (0,13) -- (26,13);
    \draw[thin] (19.5,0) -- (19.5,13); \draw[thin] (13,6.5) -- (26,6.5);
    \fill[black!30] (19.5,0) rectangle (26,6.5);
    \draw[thin] (19.5,0) rectangle (26,6.5);
    \node[lbl, anchor=north] at (13,-2) {$I_2$, and so on};
  \end{scope}
\end{tikzpicture}
\end{center}
\figcap{Halving each edge gives $2^k$ pieces; at least one inherits the
failure. Diameters fall like $2^{-n}\delta$, so the nested pieces close on a
point.}

\smallskip
Three ingredients, each traceable to Chapter 1.

\smallskip
\textbf{Something survives the nesting} --- Theorem 2.39, which rests on 2.38,
which rests on the least-upper-bound property. This is where completeness
enters, and it is the only place. Over $\Q$ the argument breaks here and
$k$-cells are not compact.

\smallskip
\textbf{The pieces shrink to nothing} --- property (c). Diameters are halved
each time.

\smallskip
\textbf{Shrinking beats any fixed radius} --- the archimedean property, which
Rudin flags explicitly in parentheses. Without it $2^{-n}\delta$ might stay
above $r$ forever.

\smallskip
Then the contradiction is immediate: $I_n$ fits inside the single set
$G_\alpha$, so one set covers it, contradicting ``no finite subcollection
covers $I_n$.'' One is finite.

\smallskip
Seen from the comparison at Definition 1.10: this proof \emph{is} the
implication ``nested intervals $\Rightarrow$ Heine--Borel'' from the ECNU
text's cycle of six equivalent completeness statements, specialised to cells.
Bisection is the standard bridge between the two.
\end{unpack}

The equivalence of (a) and (b) in the next theorem is known as the
\emph{Heine-Borel theorem}.

\ritem{2.41}{Theorem}
\emph{If a set $E$ in $\R^k$ has one of the following three properties, then it
has the other two:}
\begin{enumerate}[label=(\alph*), leftmargin=2.2em, itemsep=1pt, topsep=3pt]
  \item \emph{$E$ is closed and bounded.}
  \item \emph{$E$ is compact.}
  \item \emph{Every infinite subset of $E$ has a limit point in $E$.}
\end{enumerate}

\begin{proof}
If (a) holds, then $E \subset I$ for some $k$-cell $I$, and (b) follows from
Theorems 2.40 and 2.35. Theorem 2.37 shows that (b) implies (c). It remains to
be shown that (c) implies (a).

If $E$ is not bounded, then $E$ contains points $\mathbf{x}_n$ with
\[ |\mathbf{x}_n| > n \qquad (n = 1,2,3,\dots). \]
The set $S$ consisting of these points $\mathbf{x}_n$ is infinite and clearly
has no limit point in $\R^k$, hence has none in $E$. Thus (c) implies that $E$
is bounded.

If $E$ is not closed, then there is a point $\mathbf{x}_0 \in \R^k$ which is a
limit point of $E$ but not a point of $E$. For $n = 1,2,3,\dots$, there are
points $\mathbf{x}_n \in E$ such that $|\mathbf{x}_n-\mathbf{x}_0| < 1/n$. Let
$S$ be the set of these points $\mathbf{x}_n$. Then $S$ is infinite (otherwise
$|\mathbf{x}_n-\mathbf{x}_0|$ would have a constant positive value, for
infinitely many $n$), $S$ has $\mathbf{x}_0$ as a limit point, and $S$ has no
other limit point in $\R^k$. For if $\mathbf{y} \ne \mathbf{x}_0$, then
\begin{align*}
  |\mathbf{x}_n-\mathbf{y}| &\ge |\mathbf{x}_0-\mathbf{y}|
    - |\mathbf{x}_n-\mathbf{x}_0|\\
  &\ge |\mathbf{x}_0-\mathbf{y}| - \frac1n
    \ge \tfrac12|\mathbf{x}_0-\mathbf{y}|
\end{align*}
for all but finitely many $n$; this shows that $\mathbf{y}$ is not a limit
point of $S$ (Theorem 2.20). Thus $S$ has no limit point in $E$; hence $E$ must
be closed if (c) holds.
\end{proof}

We should remark, at this point, that (b) and (c) are equivalent in any metric
space (Exercise 2.26), but that (a) does not, in general, imply (b) and (c).
Examples are furnished by Exercise 2.16 and by the space $\mathscr{L}^2$, which
is discussed in Chap.\ 11.

\begin{center}
\begin{tikzpicture}[x=1mm, y=1mm]
  \node[rmbox, text width=25mm] (a) at (0,0)  {(a) closed and\\bounded};
  \node[rmbox, text width=24mm] (b) at (41,17) {(b) compact};
  \node[rmbox, text width=27mm] (c) at (82,0)  {(c) infinite subsets\\have limit points};
  \draw[rmarrow] (a) -- node[lbl, anchor=south east, inner sep=3pt] {2.40 + 2.35} (b);
  \draw[rmarrow] (b) -- node[lbl, anchor=south west, inner sep=3pt] {2.37} (c);
  \draw[rmarrow] (c) -- node[lbl, anchor=north, inner sep=3pt] {proved directly in 2.41} (a);
\end{tikzpicture}
\end{center}
\figcap{The cycle proved in 2.41. In a general metric space the arrow
(a)$\to$(b) is the one that breaks (Exercise 2.16, $\mathscr{L}^2$), while
(b)$\Leftrightarrow$(c) survives everywhere (2.37 and Exercise 2.26).}

\begin{deeper}[Where this sits among the completeness theorems]
Rudin's remark deserves emphasis: \textbf{(a) is the odd one out}. ``Closed and
bounded'' is not equivalent to compactness in a general metric space --- it is
a fact about $\R^k$ specifically, and the whole apparatus 2.38--2.40 exists to
establish it. Exercise 2.16 gives a closed bounded non-compact set inside $\Q$;
in $\mathscr{L}^2$ the closed unit ball fails too.

\smallskip
This is also where Chapter 2 rejoins the discussion at Definition 1.10. Of the
six equivalent formulations of completeness collected in the ECNU text, three
appear in this chapter: \textbf{Heine--Borel} is (a)$\Leftrightarrow$(b) here,
\textbf{Bolzano--Weierstrass} is 2.42, and the \textbf{nested interval}
property is 2.38. Rudin proves each separately from the least-upper-bound
property rather than deriving them from one another, so their equivalence never
becomes visible. It is worth knowing that the four statements ---
least-upper-bound, nested intervals, Heine--Borel, Bolzano--Weierstrass --- are
interderivable, and that Rudin's ordering is a choice, not a necessity.

\smallskip
Tao (\emph{Analysis II}, 1.5) takes the opposite route: he defines compactness
by sequences (every sequence has a convergent subsequence, Rudin's (c)) and
proves the open-cover version afterwards. That ordering is friendlier, and it
is why (c) is often the easier property to verify in practice --- but it does
not generalise past metric spaces, whereas 2.32 does.
\end{deeper}

\ritem{2.42}{Theorem (Weierstrass)}
\emph{Every bounded infinite subset of $\R^k$ has a limit point in $\R^k$.}

\begin{proof}
Being bounded, the set $E$ in question is a subset of a $k$-cell
$I \subset \R^k$. By Theorem 2.40, $I$ is compact, and so $E$ has a limit point
in $I$, by Theorem 2.37.
\end{proof}
